\section{ECC2K-130 arithmetic}
\label{sec:ecc2k-arith}

ECC2K-130 is the Certicom challenge
\[
  E:\ y^2 + xy = x^3 + 1 \qquad\text{over }\FF_{2^{131}},
\]
group order $4\cdot r$ with
\[
  r = 680564733841876926932320129493409985129
\]
a 129-bit prime~\cite{bailey2009,certicom1997}. Expected work is
$2^{60.9}$ iterations of a Frobenius-equivariant walk. This section
records the field, the iteration, the golden model, and the small-curve
controls. Throughput, FPGA, and the live campaign follow in
Sections~\ref{sec:ecc2k-gpu}--\ref{sec:campaign}.

\subsection{Why a type-II optimal normal basis}

$2\cdot 131+1=263$ is prime and $2$ has order $131$ modulo $263$, so
$\FF_{2^{131}}$ has a type-II optimal normal basis. With $\gamma$ a
primitive $263$rd root of unity the basis is $\beta_i=\gamma^i+\gamma^{-i}$,
$i=1,\ldots,131$, and an element is a 131-bit vector. Three
consequences, and every design in both repositories is built on all
three:\art{fpga/README.md}

\begin{center}
\small
\begin{tabular}{@{}p{0.42\textwidth}p{0.48\textwidth}@{}}
\toprule
fact & consequence \\
\midrule
$\beta_i\beta_j=\beta_{i+j}+\beta_{i-j}$
  (indices folded by $\beta_{263-k}=\beta_k$, $\beta_0=0$)
  & a product is a cyclic convolution of two symmetric 263-bit vectors \\
$\beta_i^2=\beta_{2i\bmod 263}$
  & squaring is a permutation --- free, it is wiring --- so Frobenius
    costs nothing in this basis \\
Hamming weight is basis-invariant between this index order and the
  normal-basis order
  & the iteration's $j$ and the distinguished-point test read the
    vector the hardware already holds \\
\bottomrule
\end{tabular}
\end{center}

None of that is taken on trust. The golden C model
(\texttt{fpga/model/ecc2k130.c}) builds the field independently and
checks the basis, the product rule, the squaring permutation, that
$1$ is the all-ones vector, and that the trace is the parity of the
weight. The curve is identified the same way: $E_0(\FF_2)$ has 4
points, so the Frobenius trace is $t=-1$, and the Koblitz recursion
$s_m=t\,s_{m-1}-2s_{m-2}$ gives $\#E(\FF_{2^{131}})=2^{131}+1-s_{131}
=4\cdot r$. The test then checks the two facts that make this
\emph{this} curve: $\sigma^2+\sigma+2=0$ on every point, and $n\cdot
P=\mathcal{O}$ for points multiplied by the cofactor.

\subsection{Three multiplier realisations}

A 131-bit product is realised three ways, kept bit-for-bit compatible
so a distinguished-point record re-walks under any of them.

\begin{enumerate}[leftmargin=1.2em,itemsep=0.4em]
\item \emph{Digit-serial cyclic convolution} (this library's FPGA
  core). \texttt{DIGIT} bits of the convolution per cycle; inversion
  is Itoh--Tsujii~\cite{itoh1988} along $1,2,4,8,16,32,64,65,130$,
  eight multiplications, over the core's one multiplier.
  Section~\ref{sec:fpga-core}.
\item \emph{Bitsliced software} (companion client, historical GPU
  control). Thirty-two walks in the bits of a word; the original
  packed-backend audit measured this control at $0.852294$~B/s on an
  RTX~PRO~6000.\art{crypto:ecc2k130/PACKED.md}
\item \emph{Packed polynomial-basis products} (companion client, and
  the GPU kernel imported here). Coordinates live in
  $\FF_2[c]/(m)$, $c=\gamma_1$; a 131-bit product is three $64\times
  64$ carry-less multiplies --- six \texttt{clmad} --- plus a
  reduction. Bernstein--Lange conversions to and from the normal
  basis (where the weight, the selection and the distinguished-point
  test live) are $O(K\log K)$ shift-and-mask networks. A single-product
  variant with word-parallel conversions at the boundaries is the
  first rung of the GPU ledger in Section~\ref{sec:gpu-ledger}.
\end{enumerate}

The companion also has a pipelined Karatsuba multiplier in VHDL for
the F2 engine (Section~\ref{sec:fpga-f2}).

\subsection{Two iteration functions}

The FPGA core and the historical campaign walk the Bailey et
al.~\cite{bailey2009} map
\begin{align*}
  j &= ((\HW(x_R)\gg 1)\bmod 8)+3, \qquad j\in[3,10],\\
  R' &= \sigma^j(R)+R, \qquad \sigma(x,y)=(x^2,y^2),\\
  \text{distinguished} &\iff \HW(x)\le 34
\end{align*}
(campaign cutoff $32$ is the live-fleet predicate;
Section~\ref{sec:campaign}). One step is one affine addition: one
inversion, two multiplications and a free Frobenius, since $\sigma^j$
is wiring. The inversion is eight of about eleven multiplications, so
it is $70\%$ of the step --- the single most important number in the
FPGA design.\art{fpga/README.md, ``The walk''}

Only $x$ is reported. On a Koblitz curve $-(x,y)=(x,x+y)$, so $x$ is
already the negation-class invariant; the Frobenius class is
normalised host-side. No coefficients are tracked: a walk is identified
by its seed and a collision is resolved by replaying both walks.

The packed GPU client imported into this library walks a different
function, an $r$-adding walk on the classes of $\langle\sigma,-1\rangle$
(262 points each):
\begin{align*}
  h(R) &= (\HW(x)/2)\bmod 8,\\
  k(R) &= \bigl(\textstyle\sum_e L(e)\,x_e\bigr)\cdot\HW(x)^{-1}\bmod 131,\\
  \varepsilon(R) &= \text{bit $p(R)$ of $y$},\\
  R' &= R + (-1)^{\varepsilon}\,\sigma^{k}(T_{h}),
\end{align*}
with $T_h=[a_h]P+[b_h]Q$ and a $131\times 8$-point table ($48.7$~KB)
in shared memory. $k$ shifts by one under $\sigma$ and $\varepsilon$
flips under negation, so the walk descends to classes with no
canonical representative in the loop. Fruitless 2- and 4-cycles are
avoided by advancing $h$ against the last four step tags. The two
walks share the field, the model and the 32-byte record; they do not
share a corpus.\art{ecc2k130/README.md, ``The walk''}

The exceptional case $\sigma^j(R)=\pm R$ makes the denominator zero.
The FPGA core raises \texttt{exc} and stops; dividing would produce a
point that is not on the curve, which the core would then report as
distinguished.

\subsection{Golden model and CPU walker}

\texttt{fpga/model/ecc2k130.c} is the oracle. The FPGA testbenches
compare against vectors it emits; the GPU host
(\texttt{ecc2k130/src/hostcheck.cpp}) re-slices packed words into the
same permuted type-II basis and checks equality, not equivalence up
to a basis change. \texttt{make} in \texttt{ecc2k130/} runs $1{,}904$
checks (gcc and clang in CI) covering field arithmetic, Frobenius
networks, basis conversions, table selection and the walk, with no
CUDA required.\art{ecc2k130/README.md}

The companion client's source compiles for CUDA (32-bit lanes) and
for the CPU (64-bit lanes), so the exact code that would run on a GPU
is what the CPU test suite exercises. That CPU walker is the
verification path for every GPU row below; it is not a throughput
claim.

\subsection{Small-curve controls, including ECC2K-95}

ECC2K-95 is the retired Certicom Koblitz challenge over
$\FF_{2^{97}}$, group order $4\cdot r$ with $r$ a 95-bit prime, about
$2^{44}$ iterations, solved by Harley and collaborators in
1998~\cite{harley1998}. $2\cdot 97+1=195$ is composite, so there is
no type-II optimal normal basis; the companion runs a
polynomial-basis backend with the orbit-invariant weight taken
through a linear map. The published secret
$37837308472231540269443981458$ is the check. Harley's own
measurement was $2.16\cdot 10^{13}$ iterations against a repository
prediction of $1.79\cdot 10^{13}$. The ECC2K-95 CUDA kernel itself
has run only through CPU emulation and is not a headline
throughput.\art{crypto:ecc2k130/README.md}\art{crypto:gpu/ecc2k/README.md}

Toy Koblitz curves with 21-bit and 40-bit subgroups (Section~\ref{sec:crossval})
close the pipeline to a verified logarithm in seconds. They are the
control that the iteration commutes with the class action, not a
substitute for the 129-bit group.
